\documentclass[11pt]{article}
\usepackage{mathblatt}
\begin{document}
\weit
\typeout{PROBE-BEGIN brueche-dezimalzahlen-e1-k1-s6-v17}
\begin{aufgabe}{\detokenize{brueche-dezimalzahlen-e1-k1-s6-v17}}
\begin{teile}
\teil Probe

P \streifen[0]{20}{}{} \quad Q \kreissektor{60}{} \quad R \bruchrechteck{1}{4} \quad S \bruchkreis{1}{3}
\end{teile}
\end{aufgabe}
\typeout{PROBE-END brueche-dezimalzahlen-e1-k1-s6-v17}
\typeout{PROBE-BEGIN kreis-e3-k1-s0-v1}
\begin{aufgabe}{\detokenize{kreis-e3-k1-s0-v1}}
\begin{teile}
\teil Probe

\kreissektor{180}{}
\end{teile}
\end{aufgabe}
\typeout{PROBE-END kreis-e3-k1-s0-v1}
\typeout{PROBE-BEGIN wahrscheinlichkeit-e2-k6-s3-v3}
\begin{aufgabe}{\detokenize{wahrscheinlichkeit-e2-k6-s3-v3}}
\begin{teile}
\teil Probe

\kreissektor{270}{270°}
\end{teile}
\end{aufgabe}
\typeout{PROBE-END wahrscheinlichkeit-e2-k6-s3-v3}
\typeout{PROBE-BEGIN brueche-dezimalzahlen-e1-k1-s6-v18}
\begin{aufgabe}{\detokenize{brueche-dezimalzahlen-e1-k1-s6-v18}}
\begin{teile}
\teil Probe

P \bruchkreis{1}{4} \quad Q \streifen[0]{30}{}{} \quad R \bruchrechteck{1}{6} \quad S \kreissektor{72}{}
\end{teile}
\end{aufgabe}
\typeout{PROBE-END brueche-dezimalzahlen-e1-k1-s6-v18}
\typeout{PROBE-BEGIN kreis-e3-k1-s0-v2}
\begin{aufgabe}{\detokenize{kreis-e3-k1-s0-v2}}
\begin{teile}
\teil Probe

\kreissektor{90}{}
\end{teile}
\end{aufgabe}
\typeout{PROBE-END kreis-e3-k1-s0-v2}
\end{document}
